\section{Discussion}
\subsection{What the Calculus Adds}

The central practical object is the finite coalition remainder. Suppose an
analyst plans to change a converter set point, retune a controller, and alter a
nearby network element. Standard first-order screening estimates each
increment around one baseline and then sums. The present calculation instead
asks: after every subset has moved to its own ac equilibrium and its DAE has
been rebuilt, what part of the full-portfolio characteristic change cannot be
assigned to lower-order subsets?

That question is useful even when the answer is not operationally material.
It distinguishes three failure modes of a simpler approximation:
\begin{enumerate}
    \item \emph{finite-action error}: a local Taylor model is inaccurate over
    the specified amplitude;
    \item \emph{re-equilibration error}: the operating point and algebraic
    Jacobian move jointly with the actions;
    \item \emph{interaction-order error}: pair or triple terms remain after
    all lower-order terms are removed.
\end{enumerate}
The A7--A8 result contains all three. At OP00, a frozen-affine representation
misses about 84.1\% of the finite coefficient, while exact factor anatomy
shows that its largest contribution is algebraic re-equilibration. A3--A6,
however, has a reproducible pole-only coefficient. The method thus prevents
the blanket conclusion that a characteristic externality must be merely a
power-flow artifact or a dangerous control mode.

The connected reconstruction also supplies a stringent internal check. Direct
M\"obius differencing and integrated connected traces have different numerical
error paths. Their locked-holdout agreement makes it difficult for a sign
convention, missing partition, or omitted curvature term to survive
undetected. Exact Schur closure adds a second independent identity at the
descriptor-factor level.

\subsection{Finite Portfolio Information Absent from a Baseline First-Order Screen}

A first eigenvalue sensitivity can rank actions for one mode near one
equilibrium. Higher-order and equilibrium-point sensitivities extend that
local picture \cite{nam2000eigensensitivity,zamoracardenas2013multiparameter}.
Selective modal analysis can reduce the calculation to relevant state groups
\cite{perezarriaga1982sma}. These remain the appropriate tools when the
decision question is local and the limiting modal family is known.

The finite calculus contributes when the decision variable is a named
portfolio at a non-infinitesimal amplitude. It does not approximate the
M\"obius coefficient from a derivative at zero; it evaluates the coefficient
directly and then reconstructs it by integrating the derivative over the
whole hypercube. This is why mixed re-equilibration curvature matters. The
6.3-fold A7--A8 discrepancy would not be repaired merely by better tracking of
the baseline eigenvalue.

Conversely, the critical-mode audit shows what the spectral functional can
miss if used alone: relevance to the limiting mode. The functional found
converter-control interactions reproducibly, and the pole factor verified
that some were dynamic. But these modes were strongly damped. A separate
minimum-damping search was required to show that the operative margin belonged
to a synchronous family and that the externality barely moved it. Modal
association tools---BMAC, participation, residues, condition numbers, and
homotopy---support that family identification; none is a causal proof.

\subsection{Power-System Relevance}

Three features make the experiment more than a matrix identity exercise.
First, actions are recognizable engineering changes: 20\% controller-bandwidth
adjustments, a 20-MW dispatch change, and a 10\% branch-reactance change.
Second, every coalition is required to satisfy the ac equations, dynamic
initialization, and smooth-path checks. Third, the strongest pair couples
dispatch and its incident network path, exactly the setting in which frozen
network or fixed-operating-point screening is most vulnerable.

The negative materiality result is also relevant. Interaction-detection
methods are often motivated by weak-grid risk, but interaction is not
synonymous with limiting instability. In this benchmark, conventional SCR
values provide useful context yet do not identify which modal family sets the
5\% screen. The critical mode is synchronous-electromechanical despite the
presence of three GFL plants and conspicuous converter-mode externalities.
Planning studies should therefore distinguish at least three questions:
\begin{enumerate}
    \item Is a portfolio non-additive in a chosen dynamic functional?
    \item Does the non-additivity arise in algebraic coordinates or finite
    poles?
    \item Does it materially change the mode or constraint that limits the
    operating decision?
\end{enumerate}
The first two can be positive while the third is negative.

\subsection{Limitations}

The action space contains eight coordinates and only one matched IEEE 39-bus
GFL testbed. The results establish reproducible existence, not prevalence
across converter vendors, control architectures, larger networks, or action
amplitudes. The three converter models share a common class. Grid-forming
controls, hybrid populations, protection and saturation events, unbalanced
dynamics, and topology-discontinuous actions remain outside scope.

The spectral functional is dimensionless and band dependent. Its sign and
magnitude are not direct measures of damping, instability probability,
energy, or economic cost. The normalized logarithmic frequency weight and
0.1--30 Hz band were frozen before the campaign, but another defensible
functional may rank coalitions differently.

The finite-to-connected identity requires smooth re-equilibrated interpolation
throughout the hypercube. It does not cover a path that crosses a power-flow
bifurcation, limiter discontinuity, state-dimension change, or singular
algebraic Jacobian. The exact Schur anatomy further assumes index-one behavior,
nonsingular $g_y$, and a nonsingular registered dynamic mass block. Contour
localization is conditional on pole-free boundaries and numerical occupancy
certificates; A3--A6 demonstrates that a pole externality may exist even when
a selected contour cannot certify its localized moment.

The materiality ladder is deliberately narrow. The 0.0025 damping-externality
threshold and 5\% screen are preregistered study rules, not universal
standards. E05C ended at a voltage boundary before finding material
composition. E05D could not start because all empty baselines were already
below the screen; it therefore leaves weak-grid materiality unresolved. No
coalition result from E05D was viewed.

Finally, no cycle or SCC localization, mechanism-targeted correction,
matched-cost negative control, nonlinear time-domain consequence, or blind EMT
transfer was performed. E02C validates the benchmark infrastructure only.
Consequently, this paper does not offer an industrial mitigation or claim that
a particular interaction graph is causal.

\subsection{Implications for Future Use}

The terminal findings suggest a disciplined workflow for other systems when
finite portfolio composition is the decision question.
First, define finite actions in physical units and solve every coalition
vertex. Second, resolve direct externalities against numerical uncertainty.
Third, use connected reconstruction as an implementation and anatomy check.
Fourth, split algebraic and pole factors exactly. Fifth, identify the actual
constraint-setting modal family independently of the externality search.
Only if the same family carries a material, reproducible finite effect should
an intervention experiment be designed.

This ordering limits unnecessary controller surgery. In the present testbed,
the descriptive audit explains why further tuning of the externality-dominant
converter modes would not answer the registered stability-margin question.
The correct terminal outcome is a boundary statement, not a new stress search.
